\begin{helper}
Fix \(\varepsilon,\eta,\varepsilon_1,\varepsilon_2,\beta\in\mathbb R\) and
\(n_0\in\mathbb N\).  Assume
\[
0<\varepsilon,\qquad 0<\eta,\qquad \eta<\varepsilon,\qquad
\beta=\frac12,
\]
and \(\noderef{ParameterHierarchy}(\eta,\varepsilon_1,\varepsilon_2,n_0)\).
Suppose the \(\eta\)-scale regularity event
\[
\noderef{TwoBiteRegularityEvent}
\bigl(k=\noderef{TwoBiteNaturalIndependenceScale}(1+\eta,n),
\eta,\varepsilon_1,\varepsilon_2,\omega\bigr)
\]
holds with high probability under the two-bite law with parameters
\[
n,\quad \noderef{TwoBiteNaturalReducedVertexCount}(n),\quad
\noderef{TwoBiteEdgeProbability}(\beta,n).
\]
Suppose also that for every \(\delta>0\) there is \(N\ge n_0\) such that for
all \(n>N\), all \(m,p\) with
\[
m=\noderef{TwoBiteNaturalReducedVertexCount}(n),\qquad
p=\noderef{TwoBiteEdgeProbability}\!\left(\frac12,n\right),
\]
and all \(I\subseteq\mathrm{Fin}(n)\) with
\[
|I|=\noderef{TwoBiteNaturalIndependenceScale}(1+\eta,n),
\]
one has
\[
\noderef{TwoBiteEventProbability}\!\left(n,m,p,\,
\{\,\omega:
I\text{ is independent in }\noderef{TwoBiteFinalGraph}(\omega)_0
\text{ and }
\noderef{TwoBiteRegularityEvent}
(k=|I|,\eta,\varepsilon_1,\varepsilon_2,\omega)\,\}\right)
\le
\delta\binom n{|I|}^{-1}.
\]
Then the final graph has independence number less than
\[
(1+\varepsilon)\sqrt{(n:\mathbb R)\log(n:\mathbb R)}
\]
with high probability under the two-bite law with edge parameter \(\beta\).
\end{helper}
\begin{proof}
Write
\[
M(n)=\noderef{TwoBiteNaturalReducedVertexCount}(n),\quad
p_\beta(n)=\noderef{TwoBiteEdgeProbability}(\beta,n),\quad
K(n)=\noderef{TwoBiteNaturalIndependenceScale}(1+\eta,n),
\]
and let
\[
\mathbb P_n(E)=
\noderef{TwoBiteEventProbability}(n,M(n),p_\beta(n),E).
\]
Because \(\beta=1/2\), this is the same law as the fixed-set estimate's law
with \(p=\noderef{TwoBiteEdgeProbability}(1/2,n)\).

Let \(R_n\) be the \(\eta\)-scale regularity event appearing in the statement.
By hypothesis, \(\mathbb P_n(R_n)\to1\).  Equivalently, for every
\(\rho>0\) there is a threshold \(N_R\) such that, for all \(n\ge N_R\),
\[
\mathbb P_n(R_n^c)<\rho/2. \tag{1}
\]

Apply \(\noderef{NaturalIndependenceScaleFitsTarget}\) to \(0<\eta<\varepsilon\).
There is \(N_{\mathrm{fit}}\) such that, for all \(n\ge N_{\mathrm{fit}}\),
\[
(K(n):\mathbb R)<
\noderef{TwoBiteIndependenceScale}(1+\varepsilon,n)
=(1+\varepsilon)\sqrt{(n:\mathbb R)\log(n:\mathbb R)}. \tag{2}
\]
Using \(\noderef{NaturalParameterApproximation}\) and increasing the threshold
if necessary, we may also assume
\[
0<K(n)\le n. \tag{3}
\]
Indeed, the natural scale is an integer ceiling of the real scale, and the
real scale grows while remaining \(o(n)\), so eventually it is positive and
strictly below \(n\).

Fix \(\rho>0\), set \(\delta=\rho/2\), and apply the fixed-set estimate with
this \(\delta\).  It gives a threshold \(N_F\) such that, for every
\(n>N_F\) and every \(K(n)\)-element set \(I\subseteq\mathrm{Fin}(n)\),
\[
\mathbb P_n\bigl(I\text{ is independent in }\noderef{TwoBiteFinalGraph}_0
\text{ and }R_n\bigr)
\le
\delta\binom n{K(n)}^{-1}. \tag{4}
\]
Here the regularity event in the fixed-set estimate is
\(\noderef{TwoBiteRegularityEvent}(k=|I|,\eta,\varepsilon_1,\varepsilon_2)\),
which is \(R_n\) because \(|I|=K(n)\).

Let \(B_n\) be the bad event that the desired independence-number conclusion
fails.  Unfolding \(\noderef{IndependenceNumberLess}\), on \(B_n\) there is a
finite independent set \(J\subseteq\mathrm{Fin}(n)\) with
\[
(1+\varepsilon)\sqrt{(n:\mathbb R)\log(n:\mathbb R)}
\le (|J|:\mathbb R). \tag{5}
\]
For \(n\ge N_{\mathrm{fit}}\), (2) and (5) imply \(K(n)\le |J|\).  Together
with \(0<K(n)\) from (3), choose a \(K(n)\)-element subset \(I\subseteq J\).
Independence is inherited by subsets, so \(I\) is independent in the same
final graph.  Thus
\[
B_n\cap R_n
\subseteq
\bigcup_{\substack{I\subseteq\mathrm{Fin}(n)\\ |I|=K(n)}}
\{I\text{ is independent in }\noderef{TwoBiteFinalGraph}_0\}\cap R_n.
\tag{6}
\]

For \(0\le K(n)\le n\), the number of \(K(n)\)-element subsets of
\(\mathrm{Fin}(n)\) is \(\binom n{K(n)}\).  Applying the finite union bound to
(6), and then applying (4) to each indexed set, gives
\[
\mathbb P_n(B_n\cap R_n)
\le
\binom n{K(n)}\,
\delta\binom n{K(n)}^{-1}
=\delta. \tag{7}
\]
The final equality uses positivity of \(\binom n{K(n)}\), which follows from
(3).

Finally \(B_n\subseteq R_n^c\cup(B_n\cap R_n)\).  Combining this inclusion
with the union bound, (1), and (7), we get
\[
\mathbb P_n(B_n)<\rho/2+\delta=\rho
\]
for all sufficiently large \(n\).  Since \(\rho>0\) was arbitrary, the
probability of the bad event tends to \(0\), equivalently the probability of
the desired \(\noderef{IndependenceNumberLess}\) event tends to \(1\).  This
is the asserted \(\noderef{WithHighProbability}\) conclusion.
\end{proof}
